\chapter{Documentation juridique}\label{ch:fm:legal-documentation}

Lorsque Lehman Brothers a déposé le bilan en septembre 2008, Metavante Corporation détenait un swap de taux portant sur 600~millions de dollars de notionnel avec
l'une de ses filiales, et le swap avait évolué en faveur de la filiale. Metavante n'a ni résilié ni payé~: elle s'est appuyée sur une clause du
contrat-cadre standard qui subordonne chaque paiement à l'absence de \omterm{def:fm:legal-documentation:eod}{cas de défaut} survenu à l'égard de l'autre partie. Un an
plus tard, en septembre 2009, le tribunal de faillite de New York a tranché à l'audience que «~rester sur le marché pendant un an, sans
prendre la moindre mesure~» était contraire à l'esprit du Bankruptcy Code, que Metavante avait renoncé à son droit de résilier en ne l'exerçant pas
rapidement, et qu'elle devait payer. En avril 2012, la Cour d'appel de Londres a lu la même clause dans l'administration européenne de Lehman et a jugé
que, selon le droit anglais, l'obligation de paiement d'une partie non défaillante reste suspendue tant que le défaut perdure. Quelques lignes d'un
document standard ont décidé qui devait quoi à qui, et deux juridictions les ont lues différemment.

\section{Le contrat-cadre et sa raison d'être}

Une société de trading opérant de gré à gré signe un contrat-cadre avec chaque contrepartie, et chaque transaction entre elles devient une confirmation placée sous celui-ci.
Le livre~2 a présenté le contrat-cadre ISDA et le GMRA du marché du repo~; le livre~6 a présenté la compensation avec déchéance du terme (close-out netting) qu'ils rendent possible. Le principe de conception est que
toutes les transactions relevant d'un même contrat-cadre forment un seul et même accord, de sorte qu'en \omterm{def:fm:legal-documentation:eod}{cas de défaut} elles soient résiliées ensemble et compensées en un seul montant, plutôt
que d'être écrémées par un administrateur qui exécute les transactions rentables et répudie le reste.

La forme imprimée du contrat est standard~; la négociation porte sur son annexe particulière, qui fixe les choix et modifie les clauses imprimées, et dans l'annexe
de collatéral. Ce que négocient les juristes et les trésoriers d'une société de trading tient en une courte liste~: quels événements permettent à chaque partie de résilier, quel montant de collatéral circule
et à quel moment, quel collatéral est éligible, et ce qui se passe le premier jour où un incident survient (la \cref{fig:fm:legal-documentation:stack}).

\begin{omfigure}
\begin{tikzpicture}[font=\footnotesize,>=stealth,box/.style={draw,rounded corners,align=center,minimum height=0.7cm,text width=3.3cm}]
\node[box,fill=omDef!15] (m) at (0,2.4) {contrat-cadre\\{\scriptsize forme imprimée}};
\node[box,fill=omDef!25] (s) at (0,1.2) {annexe particulière\\{\scriptsize choix, seuils}};
\node[box,fill=omProp!15] (c) at (-4.2,0) {confirmations\\{\scriptsize une par transaction}};
\node[box,fill=omThm!15] (a) at (4.2,0) {annexe de collatéral\\{\scriptsize seuil, MTA, actifs éligibles}};
\node[box,fill=gray!15,text width=2.8cm] (pb) at (-4.9,2.4) {convention de prime brokerage};
\node[box,fill=gray!15,text width=2.8cm] (cl) at (4.9,2.4) {conventions de compensation et de give-up};
\draw[->] (m) -- (s); \draw[->] (s) -- (c); \draw[->] (s) -- (a);
\draw[dashed] (pb) -- (m.west); \draw[dashed] (cl) -- (m);
\node[font=\scriptsize,align=center] at (0,-0.4) {un seul contrat-cadre~:\\résiliés et compensés ensemble};
\end{tikzpicture}

\omcaption{Les documents d'une relation de gré à gré~: le contrat-cadre et son annexe particulière, les confirmations qui en font partie, et
l'annexe de collatéral~; les \omterm{def:fm:legal-documentation:pba}{conventions de prime brokerage}, de compensation et de give-up sont placées à côté et peuvent être liées par une \omterm{def:fm:legal-documentation:ate}{clause de défaut croisé}. Schéma de principe.}\label{fig:fm:legal-documentation:stack}
\end{omfigure}

\section{Cas de défaut, événements de résiliation et compensation avec déchéance du terme}

\begin{definition}[Cas de défaut, événement de résiliation]\label{def:fm:legal-documentation:eod}
Un \emph{cas de défaut}\index{cas de défaut} est un événement imputable à une partie, comme le défaut de paiement ou de livraison, un manquement, l'insolvabilité ou
le défaut au titre d'autres contrats, qui permet à l'autre partie de résilier toutes les transactions relevant du contrat-cadre. Un \emph{événement de
résiliation}\index{evenement de resiliation@événement de résiliation} est un événement qui n'est imputable à aucune des parties, comme l'illégalité ou un changement de législation fiscale, et qui permet la résiliation des
transactions concernées.
\end{definition}

\begin{definition}[Événement de résiliation additionnel, seuil de déclenchement sur la VL, clause de défaut croisé]\label{def:fm:legal-documentation:ate}
Un \emph{événement de résiliation additionnel}\index{evenement de resiliation additionnel@événement de résiliation additionnel} est un \omterm{def:fm:legal-documentation:eod}{événement de résiliation} que les parties ajoutent dans l'annexe particulière. Un \emph{seuil de
déclenchement sur la VL}\index{seuil de déclenchement sur la VL} (NAV trigger) est un \omterm{def:fm:legal-documentation:eod}{événement de résiliation} additionnel visant un fonds~: une baisse de sa valeur liquidative supérieure aux pourcentages convenus sur
un, trois ou douze mois, ou passant sous un plancher. Une \emph{clause de défaut croisé}\index{clause de défaut croisé} fait du défaut d'une partie au titre d'autres
contrats d'emprunt, au-delà d'un montant seuil, un \omterm{def:fm:legal-documentation:eod}{cas de défaut} au titre de celui-ci.
\end{definition}

La section~2(a)(iii) du contrat-cadre, telle que la Cour d'appel l'a citée, soumet chaque obligation de paiement à «~la condition suspensive selon laquelle
aucun \omterm{def:fm:legal-documentation:eod}{cas de défaut} ou \omterm{def:fm:legal-documentation:eod}{cas de défaut} potentiel à l'égard de l'autre partie ne s'est produit et ne se poursuit~». La clause permet à la
partie non défaillante de cesser de payer pendant qu'elle décide s'il faut résilier. La durée pendant laquelle elle peut attendre est précisément le point sur lequel les deux juridictions divergent.

\begin{definition}[Clause de suspension des paiements]\label{def:fm:legal-documentation:flawed}
Une \emph{clause de suspension des paiements}\index{clause de suspension des paiements} (flawed asset) subordonne le droit d'une partie à recevoir un paiement à la survenance d'un événement que son propre défaut
fait disparaître, de sorte que l'actif qu'elle détient (le droit d'être payée) est vicié dès l'origine plutôt que retiré en cas d'insolvabilité~; la section~2(a)(iii) en
est l'exemple type.
\end{definition}

Le tribunal anglais a jugé que la créance sous-jacente n'est pas affectée et que seule l'obligation de paiement est suspendue, que la suspension n'a pas de date de fin et
se prolonge jusqu'à la purge du défaut, et que la clause ne heurte pas le principe interdisant de priver une masse insolvable de ses actifs. Le
tribunal américain de faillite, lisant les safe harbours du Bankruptcy Code comme protégeant le droit de résilier et de compenser plutôt que le droit d'attendre, a jugé que
attendre un an avait emporté renonciation au droit. Une société de trading face à une contrepartie défaillante dans la monnaie pour le défaillant doit donc prendre une décision, rapidement,
dont la réponse dépend du droit applicable et du forum.

\begin{dated}{2026-09}{Suspensions des droits en cas de défaut et collatéral des swaps compensés}\label{dat:fm:legal-documentation:stays}
\textbf{Règles américaines de suspension.} En vertu des 12 CFR 252.83 et 252.84, les contrats financiers qualifiés des entités bancaires américaines couvertes doivent prévoir que les
dispositions de transfert et de suspension des régimes américains de résolution spéciale produisent effet comme si les contrats étaient régis par le droit américain, et ne peuvent pas permettre l'exercice des
droits de défaut liés à une filiale de la partie directe qui entre en procédure d'insolvabilité. Les contreparties de ces banques acceptent ces
conditions dans leurs contrats ou au travers de protocoles de place. \textbf{Swaps compensés.} En vertu des 17 CFR 22.2, un futures commission merchant (FCM) doit traiter les
swaps compensés de ses clients et leur collatéral comme appartenant à ceux-ci, et séparer le collatéral.
\end{dated}

\section{Annexes de collatéral}

L'annexe de collatéral convertit l'exposition de crédit du contrat-cadre en transferts quotidiens. Le livre~6 a défini ses notions~: le seuil en deçà duquel aucun
collatéral n'est dû, le montant minimal de transfert en deçà duquel aucun transfert n'est effectué, et le montant indépendant (independent amount) versé quelle que soit l'exposition. L'annexe
ajoute une convention d'arrondi, une liste de collatéraux éligibles et la décote de chacun.

\begin{example}[Un appel de marge au titre d'une annexe]\label{ex:fm:legal-documentation:call}
Une annexe fixe un seuil de 2~millions de dollars, un montant minimal de transfert de 0.5~million de dollars et un arrondi de 0.1~million de dollars~; le créancier garanti détient 5
millions de dollars en liquidités. À une exposition de 10.37~millions de dollars, le montant de couverture vaut $10.37-2=8.37$ et l'appel de marge est $8.37-5=3.37$, arrondi à la hausse à 3.4
millions de dollars. À 6.8~millions de dollars, la différence, soit une restitution de 0.2~million de dollars, est inférieure au montant minimal de transfert~: rien ne bouge. Dans l'annexe du dealer A
du tutoriel, avec un montant indépendant de 5~millions de dollars et une décote de 2~\% sur les obligations, une exposition de 4.2~millions de dollars face à 3~millions de dollars de liquidités et 5
millions de dollars d'obligations, ce qui donne un appel de $4.2+5-3-4.9=1.3$ million de dollars.
\end{example}

\omcode{code/firm/docterms/firm_docterms.py}{25}{48}{Une annexe de collatéral sous forme de données, et son appel de marge quotidien~: seuil, montant indépendant, valeur décotée du
collatéral détenu, montant minimal de transfert et arrondi.}\label{lst:fm:legal-documentation:annex}

\section{Conventions de prime brokerage et de compensation}

\begin{definition}[Convention de prime brokerage]\label{def:fm:legal-documentation:pba}
Une \emph{convention de prime brokerage}\index{convention de prime brokerage} fixe les conditions auxquelles un prime broker détient, finance et prête des titres à un
client~: la marge (la \omterm{def:fm:treasury-and-funding:house}{marge maison} du chapitre~14 et tout lock-up), les droits de réhypothécation du courtier (livre~1) sur les actifs du client, ses droits de
liquider les positions et compenser les montants en \omterm{def:fm:legal-documentation:eod}{cas de défaut}, ainsi que les \omterm{def:fm:legal-documentation:eod}{cas de défaut} et les \omterm{def:fm:legal-documentation:eod}{événements de résiliation} du client, y compris les \omterm{def:fm:legal-documentation:ate}{seuils de déclenchement sur la VL}.
\end{definition}

Les \omterm{def:fm:legal-documentation:pba}{conventions de prime brokerage} sont le lieu où les \omterm{def:fm:legal-documentation:ate}{seuils de déclenchement sur la VL} comptent le plus, car elles se superposent au pouvoir discrétionnaire du courtier sur la marge~: un fonds en drawdown
peut déclencher un seuil au moment où ses courtiers relèvent leurs appels de marge, et un déclenchement auprès d'une contrepartie peut devenir un défaut auprès d'autres par le jeu du défaut croisé.
Les conventions de compensation avec un futures commission merchant (FCM, livre~1) et les conventions de give-up entre le courtier exécutant et le courtier compensateur complètent l'ensemble~;
leurs clauses sont plus courtes, car le règlement de marché de la chambre de compensation porte l'essentiel de la charge.

\section{Tutoriel~: le seuil de déclenchement}\label{tut:fm:legal-documentation:tut}

\begin{tutorial}
\textbf{Objectif.} Charger les termes des conventions d'un fonds, dérouler une année de VL et d'expositions, calculer les appels de collatéral quotidiens, trouver le premier jour où chaque
\omterm{def:fm:legal-documentation:ate}{seuil de déclenchement sur la VL} est touché, et procéder à la clôture de la première contrepartie qui résilie. \textbf{État final~:} le graphique de la VL avec les déclenchements
(la \cref{fig:fm:legal-documentation:nav}) et le tableau de résiliation (le \cref{tab:fm:legal-documentation:closeout}).
\begin{tutsteps}
\item \textbf{Le fonds.} La VL part de 400~millions de dollars~; un semestre calme, un drawdown de trois mois, une reprise partielle
(\texttt{fm\_docs.nav\_path}, données synthétiques).
\item \textbf{Les termes.} Deux prime brokers et deux dealers ont des \omterm{def:fm:legal-documentation:ate}{seuils de déclenchement sur la VL} de 21, 63 et 252 jours ouvrés, dont l'un avec un plancher de 280
millions de dollars~; trois dealers ont des annexes de collatéral (\texttt{fm\_docs.TRIGGERS}, \texttt{ANNEXES}~; termes illustratifs).
\item \textbf{Le moniteur.} \texttt{firm.docterms.first\_trip} trouve le premier déclenchement de chaque seuil~; \texttt{fm\_docs.calls} applique les annexes chaque jour.
\item \textbf{La résiliation.} Le premier dealer à se déclencher résilie~: \texttt{firm.docterms.close\_out} compense ses swaps, le collatéral qu'il détient et un
repo, avec et sans compensation.
\end{tutsteps}
\end{tutorial}

\begin{omfigure}
\begin{tikzpicture}
\begin{axis}[width=10.5cm,height=5.2cm,xlabel={\small jour ouvré},ylabel={\small VL (millions de \$)},xmin=0,xmax=251,ymin=220,ymax=500,
  tick label style={font=\footnotesize},grid=major,grid style={gray!20},table/col sep=comma]
\addplot[omDef,thick] table[x=day,y=nav]{figdata/desk/18-legal-documentation/nav.csv};
\draw[gray,dashed] (axis cs:0,280) -- (axis cs:251,280) node[pos=0.35,below,font=\scriptsize] {plancher du dealer A};
\draw[omProp,dotted,thick] (axis cs:137,220) -- (axis cs:137,470);
\draw[omProp,dotted,thick] (axis cs:150,220) -- (axis cs:150,470);
\draw[omProp,dotted,thick] (axis cs:177,220) -- (axis cs:177,470);
\draw[omProp,dotted,thick] (axis cs:180,220) -- (axis cs:180,470);
\node[font=\scriptsize,anchor=south east] at (axis cs:137,472) {B};
\node[font=\scriptsize,anchor=south west] at (axis cs:150,472) {PB1};
\node[font=\scriptsize,anchor=south east] at (axis cs:177,438) {PB2};
\node[font=\scriptsize,anchor=south west] at (axis cs:180,472) {A};
\end{axis}
\end{tikzpicture}

\omcaption{La VL du fonds sur l'année et le premier déclenchement du seuil sur la VL de chaque contrepartie~: le dealer B (une baisse de 12~\% en 21 jours) le jour 137,
le prime broker 1 (15~\% en 21 jours) le jour 150, le prime broker 2 (30~\% en 63 jours) le jour 177 et le dealer A (le plancher de 280~millions de dollars) le jour 180 (étiquettes B, PB1, PB2, A). Données~:
\texttt{fm\_docs.nav\_path}, \texttt{fm\_docs.trips}.}\label{fig:fm:legal-documentation:nav}
\end{omfigure}

Le fonds atteint son sommet à 447.7~millions de dollars le jour 45 et tombe à 247.7~millions de dollars le jour 204, soit 44.7~\% sous ce sommet. Les seuils se déclenchent dans l'ordre de leur
sévérité, et non selon les difficultés du fonds~: le seuil de 12~\% en un mois du dealer B se déclenche le jour 137 alors que le fonds est encore à 369.3~millions de dollars, soit 12.4~\% sous son niveau
un mois plus tôt~; les prime brokers suivent dans les six semaines~; le plancher du dealer A se déclenche en dernier, le jour 180, à 275.2~millions de dollars. Quatre contreparties peuvent
résilier dans les 43 jours ouvrés, ce qui est le risque que le trésorier du chapitre~14 doit anticiper~: les termes des documents déterminent le moment où le financement
s'arrête.

\begin{omfigure}
\begin{tikzpicture}
\begin{axis}[width=10.5cm,height=5cm,xlabel={\small jour ouvré},ylabel={\small millions de \$},xmin=0,xmax=251,ymin=-3,ymax=20,
  tick label style={font=\footnotesize},grid=major,grid style={gray!20},table/col sep=comma,
  legend cell align=left,legend style={font=\scriptsize,at={(0.5,-0.3)},anchor=north,/tikz/every even column/.append style={column sep=8pt}},legend columns=2]
\addplot[omDef,thick,const plot] table[x=day,y=held_A]{figdata/desk/18-legal-documentation/collateral.csv};
\addplot[omDef!50,thin] table[x=day,y=exp_A]{figdata/desk/18-legal-documentation/collateral.csv};
\addplot[omThm,thick,const plot] table[x=day,y=held_B]{figdata/desk/18-legal-documentation/collateral.csv};
\addplot[omThm!50,thin] table[x=day,y=exp_B]{figdata/desk/18-legal-documentation/collateral.csv};
\legend{détenu par le dealer A,exposition du dealer A,détenu par le dealer B,exposition du dealer B}
\end{axis}
\end{tikzpicture}

\omcaption{Collatéral que le fonds a remis à deux dealers, au regard de leur exposition à son égard. L'annexe du dealer A ajoute un montant indépendant de 5~millions de dollars, si bien
qu'il détient en permanence davantage que son exposition~; le seuil de 2~millions de dollars du dealer B fait qu'il détient moins. Le seuil de 10~millions de dollars du dealer C n'est jamais atteint.
Données~: \texttt{fm\_docs.calls}.}\label{fig:fm:legal-documentation:collateral}
\end{omfigure}

Sur l'année, les annexes donnent lieu à des mouvements de collatéral pendant 87 jours pour le dealer A (44 appels, 43 restitutions) et 39 pour le dealer B, et jamais pour le dealer C, dont le seuil de 10
millions de dollars est supérieur à sa plus forte exposition. Le dealer B résilie le jour 137.

\begin{table}[ht]
\centering\small
\begin{tabular}{lr}
\toprule
dealer B, jour 137 & millions de \$ \\
\midrule
swaps~: exposition au mid & 2.84 \\
swaps~: coût de remplacement ajouté par le dealer & 3.00 \\
moins le collatéral détenu par le dealer & $-$1.20 \\
solde des swaps, dû par le fonds & 4.64 \\
repo~: collatéral excédentaire détenu par le dealer, dû au fonds & $-$6.00 \\
\midrule
net avec compensation (dû au fonds) & $-$1.36 \\
sans compensation & le fonds paie 4.64~; il réclame 6.00 \\
\bottomrule
\end{tabular}
\caption{Résiliation des conventions du dealer B avec le fonds à la date de son déclenchement. Données~: \texttt{fm\_docs.close\_out\_first}.}\label{tab:fm:legal-documentation:closeout}
\end{table}

Avec une clause de compensation couvrant les conventions du dealer, la dette du fonds sur les swaps et celle du dealer sur le repo se compensent en 1.36~millions de dollars dus au
fonds. Sans clause, le fonds doit payer 4.64~millions de dollars au titre des swaps et réclamer séparément 6~millions de dollars au titre du repo, ce qui ne représente la même somme
que si le dealer paie~; l'exposition brute au crédit du dealer est de 6~millions de dollars, et non la nette.

\section{Adhésions à la bourse et give-ups}

L'adhésion à la bourse (livre~1) est elle-même un contrat~: le règlement de marché lie le membre, ses traders et ses systèmes, et la convention de compensation avec un
membre compensateur garantit les transactions du membre auprès de la chambre de compensation. Une société de trading qui exécute par l'intermédiaire de plusieurs courtiers et compense par un seul signe
des conventions de give-up (les lignes de give-up du livre~2) qui permettent aux courtiers exécutants de transmettre ses transactions à son compensateur, dans des limites que celui-ci fixe. Les
documents sont plus courts qu'un contrat-cadre et changent moins, mais les limites qu'ils contiennent sont opérationnelles~: une ligne de give-up saturée interrompt le trading.

\begin{method}[Négocier les documents d'un fonds]\label{met:fm:legal-documentation:negotiate}
\begin{enumerate}
\item Lister chaque convention et, pour chacune, ses \omterm{def:fm:legal-documentation:eod}{cas de défaut}, ses \omterm{def:fm:legal-documentation:eod}{événements de résiliation} et ses seuils de déclenchement, avec leurs niveaux et leurs délais de régularisation.
\item Simuler la VL du fonds sur une mauvaise année et trouver l'ordre dans lequel les seuils se déclenchent~; élargir celui qui est le plus serré.
\item Aligner les seuils, les montants minimaux de transfert et le collatéral éligible sur les prévisions de trésorerie (chapitre~14).
\item Négocier la compensation entre conventions avec un même groupe de contreparties, et des seuils de défaut croisé assez élevés pour éviter tout effet de cascade.
\item Conserver les termes sous forme de données, et les surveiller chaque jour.
\end{enumerate}
\end{method}

\section{Construction~: les termes des conventions sous forme de données}\label{bld:fm:legal-documentation:docterms}

\begin{build}
\bfield{Rôle} Aucune construction véritable~: le sujet est la documentation juridique. L'outil d'analyse conserve les termes des conventions sous forme de données et calcule
ce qu'ils impliquent jour par jour.
\bfield{Interface} \texttt{firm.docterms}~: \texttt{Annex}, \texttt{held\_value}, \texttt{call}~; \texttt{Trigger}, \texttt{first\_trip}~;
\texttt{close\_out}.
\bfield{Règles} Le collatéral inéligible ne compte pour rien~; aucun transfert sous le montant minimal de transfert~; les appels sont arrondis au supérieur et les restitutions à l'inférieur.
les seuils comparent la VL de chaque jour à la VL d'une fenêtre antérieure et au plancher.
\bfield{Tests d'acceptation} \texttt{code/firm/docterms/tests/}~: appels d'annexe sur des chiffres calculés à la main~; déclenchements de seuils par baisse et par plancher~; résiliation avec et
sans compensation.
\bfield{Pour aller plus loin} Délais de régularisation et délais de préavis~; seuils sur la VL en fin de mois~; une cascade de défauts croisés entre les contreparties du fonds.
\end{build}

\begin{omsources}
\item Lomas v JFB Firth Rixson Inc [2012] EWCA Civ 419.
\item Davis Polk, «~Lehman Bankruptcy Court Rules \dots~», note client, 29 septembre 2009, sur In re Lehman Brothers Holdings Inc.
(Metavante).
\item 12 CFR 252.83 et 252.84~; 17 CFR 22.2.
\end{omsources}

\section{Exercices}

\begin{exercise}[$\star$]\label{exo:fm:legal-documentation:1}
Sous l'annexe de l'\cref{ex:fm:legal-documentation:call}, quels mouvements de collatéral se produisent pour des expositions de 2.3~millions de dollars et de 10.37~millions de dollars, sans collatéral détenu~?
\end{exercise}

\begin{exercise}[$\star$]\label{exo:fm:legal-documentation:2}
Quelle est la différence entre un \omterm{def:fm:legal-documentation:eod}{cas de défaut} et un \omterm{def:fm:legal-documentation:eod}{événement de résiliation}~? Lequel est un \omterm{def:fm:legal-documentation:ate}{seuil de déclenchement sur la VL}~?
\end{exercise}

\begin{exercise}[$\star$]\label{exo:fm:legal-documentation:3}
À quels jours les quatre seuils du tutoriel se déclenchent-ils, et lequel se déclenche le premier~?
\end{exercise}

\begin{exercise}[$\star\star$]\label{exo:fm:legal-documentation:4}
Pourquoi Metavante n'a-t-elle ni payé ni résilié, et qu'a décidé la cour de New York~?
\end{exercise}

\begin{exercise}[$\star\star$]\label{exo:fm:legal-documentation:5}
Qu'a décidé la Court of Appeal anglaise au sujet de la section 2(a)(iii), et en quoi diffère-t-elle de la décision rendue à New York~?
\end{exercise}

\begin{exercise}[$\star\star$]\label{exo:fm:legal-documentation:6}
Pourquoi le dealer A détient-il davantage de collatéral que son exposition, et pourquoi le dealer C n'en détient-il aucun~?
\end{exercise}

\begin{exercise}[$\star\star\star$]\label{exo:fm:legal-documentation:7}
\emph{Programmation.} Resserrer le seuil à 63~jours du prime broker 2, de 30~\% à 20~\%, et relancer \texttt{fm\_docs.trips}. Lequel se déclenche-t-il en premier désormais~?
\end{exercise}

\begin{exercise}[$\star\star\star$]\label{exo:fm:legal-documentation:8}
\emph{Trouver la faille.} «~Nos \omterm{def:fm:legal-documentation:ate}{seuils de déclenchement sur la VL} sont généreux~: aucun ne se déclenche tant que nous ne perdons pas 30~\%.~»
\end{exercise}

\section{Problème~: le seuil de déclenchement}

\begin{problem}[{Problème du week-end --- le seuil de déclenchement}]\label{pb:fm:legal-documentation:1}
Le general counsel d'un fonds doit revoir ses conventions avec deux prime brokers et trois dealers avant une année volatile.

\medskip\noindent\textbf{Partie~I --- Les documents.}
\begin{enumerate}
\item Pourquoi un contrat-cadre constitue-t-il un accord unique, et que cela apporte-t-il à la partie non défaillante~?
\item Définir un \omterm{def:fm:legal-documentation:eod}{cas de défaut}, un \omterm{def:fm:legal-documentation:eod}{événement de résiliation}, un \omterm{def:fm:legal-documentation:ate}{événement de résiliation additionnel} et un \omterm{def:fm:legal-documentation:ate}{seuil de déclenchement sur la VL}.
\item Définir une \omterm{def:fm:legal-documentation:ate}{clause de défaut croisé} et une \omterm{def:fm:legal-documentation:flawed}{clause de suspension des paiements}.
\item Que dit la section 2(a)(iii), telle que la Court of Appeal l'a citée~?
\end{enumerate}

\medskip\noindent\textbf{Partie~II --- Les affaires.}
\begin{enumerate}[resume]
\item Résumer l'affaire Metavante et son issue.
\item Résumer l'affaire Lomas et la solution retenue sur la suspension.
\item Que doit décider une partie non défaillante après le défaut d'une contrepartie, et dans quel délai~?
\item Que changent les règles américaines de stay pour les contreparties des grandes banques américaines~?
\end{enumerate}

\medskip\noindent\textbf{Partie~III --- L'année.}
\begin{enumerate}[resume]
\item Donner le sommet, le creux et le drawdown de la VL.
\item Donner le premier déclenchement de chaque seuil et sa raison.
\item Donner les mouvements de collatéral sous chaque annexe, et expliquer le dealer C.
\item Calculer l'appel de marge de l'\cref{ex:fm:legal-documentation:call}.
\item Donner le montant de résiliation du dealer B avec et sans set-off.
\end{enumerate}

\medskip\noindent\textbf{Partie~IV --- L'examen.}
\begin{enumerate}[resume]
\item Définir une \omterm{def:fm:legal-documentation:pba}{convention de prime brokerage} et citer les termes qui interagissent avec la \omterm{def:fm:treasury-and-funding:house}{marge maison} du chapitre 14.
\item Quels seuils faudrait-il renégocier en premier, et pourquoi~?
\item Que fait le défaut croisé à l'ordre des déclenchements~?
\item Pourquoi le set-off compte-t-il davantage lorsque la contrepartie est fragile~?
\item Qu'apportent les conventions de give-up à l'ensemble~?
\item Énoncer le \emph{résultat nommé}~: le premier jour où une VL en baisse déclenche un \omterm{def:fm:legal-documentation:ate}{événement de résiliation additionnel} selon les termes de chaque contrepartie, et le
montant de résiliation après set-off.
\item En deux phrases, rédiger la recommandation de l'examen.
\end{enumerate}
\end{problem}

\section{Questions d'entretien}

\begin{interviewq}[$\star$ \iqroles{trader}]\label{iq:fm:legal-documentation:1}
Qu'est-ce qu'un montant minimal de transfert, et pourquoi une annexe en comporte-t-elle un~?
\end{interviewq}

\begin{interviewq}[$\star$ \iqroles{risque}]\label{iq:fm:legal-documentation:2}
Qu'est-ce qu'un \omterm{def:fm:legal-documentation:ate}{seuil de déclenchement sur la VL}, et pourquoi les contreparties en demandent-elles aux fonds~?
\end{interviewq}

\begin{interviewq}[$\star\star$ \iqroles{développeur}]\label{iq:fm:legal-documentation:3}
Implémentez un appel de marge avec un seuil, un montant indépendant, des décotes, un montant minimal de transfert et un arrondi. Quels sont les cas limites~?
\end{interviewq}

\begin{interviewq}[$\star\star$ \iqroles{risque}]\label{iq:fm:legal-documentation:4}
Votre contrepartie a fait défaut et vous lui devez de l'argent sur les swaps. Quelles sont vos options~?
\end{interviewq}

\begin{interviewq}[$\star\star$ \iqroles{trader, risque}]\label{iq:fm:legal-documentation:5}
Pourquoi la compensation avec déchéance du terme compte-t-elle pour le capital qu'une banque détient face à une contrepartie~?
\end{interviewq}

\begin{interviewq}[$\star\star\star$ \iqroles{risque, chercheur}]\label{iq:fm:legal-documentation:6}
Comment estimeriez-vous la probabilité qu'un fonds déclenche au moins un seuil sur la VL au cours de l'année suivante~?
\end{interviewq}
